\chapter{数的表示与算术电路}

硬件只保存 bit。所谓整数、正负号、溢出，都是对 bit 模式的解释规则。若不先讲清楚表示，后面的加法器、减法器、比较器和 ALU 就会像在做抽象数学；事实上它们处理的始终是固定宽度的一组电线。

\section{一、无符号数就是按权重求和}

4-bit 无符号数 \code{b3 b2 b1 b0} 的值是：

\begin{lstlisting}
value = b3*8 + b2*4 + b1*2 + b0
\end{lstlisting}

因此 \code{1011} 表示 11，因为它等于 8+0+2+1。无符号数没有负数概念，所有 bit 都贡献非负权重。4-bit 能表达的范围是 0 到 15，8-bit 能表达的范围是 0 到 255。

\begin{pdfdiagram}[x=1cm,y=1cm]
  \node[hw label,text=BookNavy] at (4.2,2.25) {4-bit 无符号权重};
  \foreach \x/\bit/\w in {0/1/8,1/0/4,2/1/2,3/1/1} {
    \node[hw state,minimum width=1.05cm] (b\x) at (1.6+\x*1.75,1.35) {\bit};
    \node[hw label] at (1.6+\x*1.75,0.62) {权重 \w};
  }
  \draw[BookAccent,line width=0.6pt,-{Latex[length=2mm,width=1.5mm]}] (1.6,0.2) -- (1.6,-0.25);
  \draw[BookAccent,line width=0.6pt,-{Latex[length=2mm,width=1.5mm]}] (5.1,0.2) -- (5.1,-0.25);
  \draw[BookAccent,line width=0.6pt,-{Latex[length=2mm,width=1.5mm]}] (6.85,0.2) -- (6.85,-0.25);
  \node[hw label,align=center] at (4.2,-0.65) {1011 = 1*8 + 0*4 + 1*2 + 1*1 = 11};
\end{pdfdiagram}
\diagramnote{无符号数的每一位都有固定非负权重，硬件只保存 bit，数值来自解释规则。}

位宽是硬件限制，不是显示格式。一个 4-bit 加法器只能输出 4 个结果 bit，再加上可能的进位。若 15 加 1，低 4 位回到 \code{0000}，最高位外产生一个进位。这个进位不是“可有可无的小细节”，它是无符号运算是否超出范围的证据。

\section{二、补码让负数也能走加法器}

硬件希望用同一个加法器处理加法和减法。补码表示提供了这个能力。以 4-bit 为例，最高位不再只是 8，而是解释成负权重 -8：

\begin{lstlisting}
value = b3*(-8) + b2*4 + b1*2 + b0
\end{lstlisting}

于是 \code{1111} 的值是 -8+4+2+1=-1，\code{1110} 的值是 -8+4+2+0=-2，\code{1000} 的值是 -8。4-bit 补码范围是 -8 到 7。

求一个数的相反数，可以按“取反加一”理解：

\begin{lstlisting}
x = 0011  // 3
~x = 1100
~x + 1 = 1101  // -3 in 4-bit two's complement
\end{lstlisting}

为什么这个规则适合硬件？因为取反可以用一排 NOT 门完成，加一可以交给加法器。于是 a-b 可以变成：

\begin{lstlisting}
a - b = a + (NOT b) + 1
\end{lstlisting}

这样，减法器不必完全另做一套。只要在 b 输入前加反相控制，并把最低位进位输入设为 1，同一个加法器就能执行减法。

\begin{pdfdiagram}[x=1cm,y=1cm]
  \node[hw state,minimum width=1.2cm] (a) at (0.8,1.75) {a};
  \node[hw state,minimum width=1.2cm] (b) at (0.8,0.6) {b};
  \node[hw control,minimum width=1.2cm] (sub) at (0.8,-0.45) {sub};
  \node[hw compute,minimum width=1.5cm] (xor) at (3.1,0.6) {逐位 XOR};
  \node[hw compute,minimum width=1.65cm] (adder) at (5.8,1.15) {加法器};
  \node[hw state,minimum width=1.4cm] (res) at (8.35,1.15) {result};
  \draw[hw bus] (a.east) -- (adder.west);
  \draw[hw bus] (b.east) -- (xor.west);
  \draw[hw bus] (sub.east) -- (xor.west);
  \draw[hw bus] (xor.east) -- node[above,hw label] {b 或 NOT b} (adder.west);
  \draw[hw bus] (sub.east) .. controls (3.1,-0.85) and (5.0,-0.45) .. node[below,hw label] {carry\_in} (adder.south);
  \draw[hw bus] (adder.east) -- (res.west);
  \node[hw label,align=center] at (4.7,-1.2) {sub=0：a+b；sub=1：a+(NOT b)+1，也就是 a-b};
\end{pdfdiagram}
\diagramnote{补码减法把“取反”和“加一”变成控制信号，让减法复用加法器。}

\section{三、同一组 bit 会有不同溢出解释}

同一组 bit 可以按无符号解释，也可以按补码解释。因此溢出判断也不同。

\begin{longtable}{p{0.2\linewidth}p{0.34\linewidth}p{0.34\linewidth}}
\toprule
解释方式 & 关注点 & 例子 \\
\midrule
无符号 & 是否产生最高位外的进位 & 15 + 1 超出 4-bit 无符号范围 \\
补码有符号 & 正正得负或负负得正 & 7 + 1 得到 \code{1000}，被解释为 -8 \\
\bottomrule
\end{longtable}

看一个具体例子。4-bit 下：

\begin{lstlisting}
0111 + 0001 = 1000
\end{lstlisting}

若按无符号解释，这是 7+1=8，没有问题。若按补码解释，\code{0111} 是 7，\code{0001} 是 1，结果 \code{1000} 是 -8，正数加正数却得到负数，所以发生有符号溢出。

硬件不会替使用者决定哪种解释才是“正确数学”。它只提供证据：结果 bit、进位、符号位变化、溢出标志。解释规则由控制信号和上层约定决定。

\section{四、移位器也是算术部件}

左移一位通常相当于乘以 2，但前提是移出去的 bit 没有携带需要保留的信息。右移更容易出差别。逻辑右移在高位补 0；算术右移保留符号位，用于补码有符号数。

\begin{lstlisting}
logical right shift:    1000 >> 1 = 0100
arithmetic right shift: 1000 >> 1 = 1100
\end{lstlisting}

若 \code{1000} 按 4-bit 补码解释为 -8，算术右移得到 \code{1100}，也就是 -4；逻辑右移得到 \code{0100}，也就是 4。两个结果差异巨大。原因不是移位器“出错”，而是控制信号选择了不同硬件规则。

硬件中的移位器可以很简单，一次只移一位；也可以做成桶形移位器，让多级选择器按移位量一次选择结果。简单移位器面积小但可能需要多个周期，桶形移位器速度快但组合逻辑更多。

\section{五、比较可以来自减法和标志位}

相等检测可以用 XOR 或 XNOR 网络。若每一位都相同，整体相等。大小比较常常可以借助减法结果和标志位。若 a-b 的结果为 0，则 a 等于 b；若按无符号解释，可以看借位或进位；若按补码解释，还要结合符号位和溢出。

这再次说明，算术电路不只是“算出一个数”。它还要产生后续控制需要的证据。ALU 会把加法、减法、按位逻辑、移位和比较集中起来，同时输出 Zero、Carry、Overflow、Negative 等标志位。下一章就把这些能力收进一个统一组合部件。

\begin{classicnote}
本章的核心标注是“同一组 bit，多种解释”。硬件只搬运和变换 bit，解释由编码规则和控制信号决定。

\begin{itemize}
\item 纸上实验：用 4-bit 同时解释 \code{1000}：无符号是 8，补码是 -8；再计算 \code{0111 + 0001}，分别判断无符号进位和有符号溢出。
\item 位级证据：补码取负可写成按位取反再加一；减法可复用加法器执行 \code{a + (~b) + 1}。
\item 标志证据：Carry 适合无符号边界，Overflow 适合补码有符号边界；不能把两者混成同一个“溢出”。
\item 检查题：\code{1111 + 0001 = 0000}。它在无符号和补码解释下分别意味着什么？哪些标志应当被置位？
\end{itemize}

经典教材常把整数表示放在机器级程序前面讲，是因为后面的比较、分支、地址计算和溢出 bug 都来自这套位级解释。先学会看 bit，再谈“数”。
\end{classicnote}
